\chapter{Einleitung}

	Ziel unserer Projektarbeit war die Gestaltung sowie Implementierung eines Computerspiels im \textit{Pixel Art}-Stil mit einem \textit{Point-and-Click}-Interaktionsmodell. \\
    Die Idee für die Art des Spiels resultierte aus einer gemeinsamen Erarbeitung. Zunächst entstanden zehn sehr unterschiedliche Spielideen, welche
    sich jeweils aus Ideen aller zehn Gruppenmitglieder zusammensetzten. Schließlich wurde sich mittelns einer
    Abstimming für die Spielidee entschieden, welche sich dann zum \textit{CZGame} entwickelte.
    
	Aufgrund der Komplexität des Themas wurde die Arbeit auf mehrere Gruppen aufgeteilt, welche sich jeweils
	auf einen Teilbereich des Spiels fokussierten, jedoch auch untereinander eng zusammenarbeiteten.
	Es wurde sich für folgende Gruppen bzw. Teilbereiche entschieden:
	
	\begin{itemize}
		
		\item Die Gruppe \textit{\textbf{Grafikdesign}}, bestehend aus \textbf{Antonia}, \textbf{Jonas G.} und
			\textbf{Gustav}, beschäftigte sich damit, die Hintergrundgrafiken des Spiels zu gestalten.
	
		\item In der Gruppe \textit{\textbf{Kampfkonzept}} beschäftigten sich \textbf{Sophie} und \textbf{Emil}
			mit der Gestaltung von im Spiel stattfindenden Duellen zwischen der Spielfigur und Lehrkräften
			der Schule.
	
		\item \textbf{Jonas L.} und \textbf{Rita} beschäftigten sich in der Gruppe \textit{\textbf{Minigames}} mit dem
			Entwerfen kleiner spielerischer Herausforderungen. Sie überlegten  dazu zu verschiedenen
			Schulfächern passende Minipielkonzepte.
	
		\item \textbf{Wu-Yi} und \textbf{Arthur} beschäftigten sich in der Gruppe \textit{\textbf{Item- \& Charakterdesign}}
			mit dem Entwerfen der Grafiken zur Darstellung verschiedener Lehrpersonen sowie der Spielfigur und der Items.
	
		\item \textbf{Ashley} komponierte \textit{\textbf{Musik}}, die während des Spielens im Hintergrund laufen sollte.
	
	\end{itemize}

	Alle Gruppen sollten zudem ihre Konzepte und Gestaltungen schlussendlich auch entsprechend in das Spiel
	einbinden bzw. implementieren. Dabei erhielt jede Gruppe eine sehr große Unterstützung von Ashley, die
	ihr sehr großes Vorwissen bei der Programmierung u.A. in Form der Implementierung der Spiel-Engine einbrachte
	und sehr viel zusätzliche Zeit dafür investierte.
